\documentclass[12pt]{article}
\usepackage{graphicx} % Required for inserting images
\usepackage[margin=0.5in]{geometry}
\usepackage{amsmath, amssymb}
\usepackage{mathrsfs}


\usepackage{soul}


% Dr. David Brown Commands
\newcommand{\ds}{\displaystyle}
\newcommand{\R}{\mathbb{R}}
\newcommand{\N}{\mathbb{N}}
\newcommand{\Q}{\mathbb{Q}}
\newcommand{\C}{\mathbb{C}}


% Commands to get script r
\usepackage{calligra}
\DeclareMathAlphabet{\mathcalligra}{T1}{calligra}{m}{n}
\DeclareFontShape{T1}{calligra}{m}{n}{<->s*[2.2]callig15}{}
\newcommand{\scripty}[1]{\ensuremath{\mathcalligra{#1}}}

% Unit commands
\newcommand{\degree}{\text{°}}
\newcommand{\unitSpeed}{\frac{\text{m}}{\text{s}}}
\newcommand{\unitAcceleration}{\frac{\text{m}}{\text{s}^2}}

% Constant commands
\newcommand{\avogadro}{6.022\times10^{23}}

\newcommand{\boltzmannConstK}{1.381\times10^{-23}\frac{\text{J}}{\text{K}}}
\newcommand{\boltzmannConstInEV}{8.617\times10^{-5}\frac{\text{eV}}{\text{K}}}
\newcommand{\planckConst}{6.626\times10^{-34}\text{J}\cdot\text{s}}

\newcommand{\atmP}{1.01325\times10^5\,\text{Pa}}

% Known Value Commands
\newcommand{\electronMass}{9.109\times10^{-31}\,\text{kg}}
\newcommand{\protonMass}{1.673\times10^{-27}\,\text{kg}}

% Commands to easily write partial derivatives
\newcommand{\partDeriv}[2]{\frac{\partial #1}{\partial #2}}
\newcommand{\doublePartDeriv}[2]{\frac{\partial^2 #1}{\partial^2 #2}}


\title{Problem 7.51}
\author{Gabriel Decker}


\begin{document}


\maketitle
\begin{flushleft}



In this problem, we will consider a lightbulb's tungsten filament, treating it as a blackbody.
We're given that it runs at $T=3000\,\text{K}$ and that it has an emissivity of $e=1/3$.
In this problem, we will use the blackbody intensity $P/A$ equation and the equation for the photon energy density (spectrum) per photon energy.
These are the following equations.
\begin{equation}
    P/A=\sigma eT^4
\end{equation}
    with $\sigma=5.67\times10^{-8}\,\text{W/m}^2\text{K}^4$.
\begin{equation}
    U/V=\frac{8\pi(kT)^3}{(hc)^3}\int_0^\infty\frac{x^3}{e^{x}-1}dx
    \label{eq:energy-density}
\end{equation}
with $x=\epsilon/kT$.
Also useful is the photon wavelength to energy conversion.
\begin{equation}
    \epsilon=\frac{hc}{\lambda}
    \label{eq:energy-from-wavelength}
\end{equation}
\begin{equation}
    \lambda=\frac{hc}{\epsilon}
    \label{eq:wavelength-from-energy}
\end{equation}\\~\\~\\


\textbf{Part a}\\
First, we're given that the power of the bulb is $P=100\,\text{W}$.
With this information, we can determine the surface area $A$ of the bulb.
$$P/A=\sigma eT^4$$
$$\implies A=\frac{P}{\sigma eT^4}$$
$$\implies A=\frac{100\,\text{W}}{(5.67\times10^{-8}\,\text{W/m}^2\text{K}^4)(1/3)(3000\,\text{K})^4}$$
$$\implies A=6.5321\times10^{-5}\,\text{m}^2$$\\~\\~\\



\textbf{Part b}\\
We know from page 292 of the text that the planck spectrum peaks at $x=2.82$ or at energy $\epsilon=2.82\cdot kT$.
This implies that the peak photon energy in the spectrum for $T=3000\,\text{K}$ is $\epsilon=1.16833\times10^{-19}\,\text{J}$.
Plugging this into equation \ref{eq:wavelength-from-energy}, we get $\lambda=1.700\times10^{-6}\,\text{m}$, which is firmly in the thermal range.\\~\\~\\



\textbf{Part c}\\
I generated the following plot using python.
The shaded region is energy density due to the photons with wavelength in the visible spectrum.
Clearly, most of the energy is spent creating heat.\\
\includegraphics[width=0.7\linewidth]{./energy_vs_spectrum.png}\\~\\~\\



\textbf{Part d}\\
Now, we'll calculate the fraction of the bulb's total photon energy output (because the blackbody spectrum is the same as the planck spectrum) that is due to photons in the visible spectrum.
The total energy density of the spectrum is the following.
$$\frac{U}{V}=\frac{8\pi^5(kT)^4}{15(hc)^3}$$
The contribution due to just the visible spectrum can be determined by adjusting the bounds in equation \ref{eq:energy-density}.
This was done using python.
It showed that the fraction of energy coming from the visible range is about 8.1\%.
This corroborates the story shown by the plot.\\~\\~\\



\textbf{Part e}\\
To get more efficient (more energy going into visible light), we must increase the temperature.
This can be shown by the peak energy equation $\epsilon=2.82\cdot kT$.
Clearly, we want more of the peak of the spectrum to be covered, so we want the position of the peak to move towards the visible energies, which are higher than what it is at at $T=3000\,\text{K}$.
Since the peak energy equation shows that it is proportional to $T$, this implies that increasing $T$ increases the peak energy.\\~\\~\\



\textbf{Part f}\\
The theoretically most efficient blackbody lightbulb can be determined by finding the temperature that has the greatest fraction of visible light energy.
I used python to do this by simply computing and comparing the results for many temperatures, and it resulted in the temperature $T=7043\,\text{K}$ being having the highest fraction of 39.3\%.



\end{flushleft}

\end{document}
